\documentclass[a4paper]{article}
\usepackage{luatexja}
\usepackage{amsmath}
\usepackage{amssymb}
\usepackage{geometry}
\usepackage{hyperref}

\geometry{margin=25mm}

\title{金属における常温核融合：\\
正八面体の籠にある戻す席}

\author{https://github.com/hypernumbernet}
\date{\today}

\begin{document}
\maketitle

\begin{abstract}
常温核融合とは、この理論では、室温にある同符号の対の遅い荷電の入口である。一軸に沿ってそれを引き受ける相手は、同符号のイオンの籠である。正イオンの正八面体は、中心にさらに正のイオンを一つ保持し、電子をその中心から押し出す。金属の電子は遮蔽ではない。同じ正八面体の一軸上の二つの陽子は、籠の半径が \(2\ell\) であるとき、間隔を増す成分が、間隔 \(2\ell\) において、イオン電荷が少なくとも一ならば厳密に負である。第一の電荷の節のすぐ外側では、成分は \(+\infty\) へ向かう。相手の斥力が発散し、六つのイオンは節から離れたままである。節と間隔 \(2\ell\) とのあいだで成分は厳密に減り、ちょうど一つの間隔で消える。節に向かう側では成分は外向きであり、間隔 \(2\ell\) に向かう側では内向きであり、どちらの側も対をその間隔へ戻す。籠は対を節の向こうへ運ばない。席は二つの長さの比を定めず、反応率も、熱の過剰も、実験のあいだの選択も導出しない。
\end{abstract}

\section{金属に問われた入口}
\label{sec:reports}

フライシュマンとポンスは、重水の電気分解によってパラジウムに重水素を置き、中性子およびトリチウムとともに、熱の過剰を報告した~\cite{fleischmann1989}。同じ年に、ジョーンズらは、電解セルからの低い収量の中性子を報告した~\cite{jones1989}。

日本では、山口と西岡が、パラジウムの板に重水素を吸蔵させ、一つの面に酸化物の障壁を置き、その板を真空に保った。彼らは、中性子のバーストと、ガスの放出と、熱の過剰とを報告した~\cite{yamaguchi1990}。荒田と張は、閉じたパラジウム陰極の内側に微細なパラジウムを置き、重水を電気分解し、熱の過剰とヘリウム4とを報告した~\cite{arata1999}。水野らは、パラジウム陰極の上に析出した物質の同位体の異常を報告した~\cite{mizuno1996}。岩村、坂野、伊藤は、\(343\,\mathrm{K}\)、一気圧の重水素ガスを、酸化カルシウムを含むパラジウムの多層膜に透過させた。彼らは、表面のセシウムが減ってプラセオジムが現れ、ストロンチウムが減ってモリブデンが現れたと報告した~\cite{iwamura2002}。新エネルギー・産業技術総合開発機構の新水素エネルギー実証は、1998年に閉じ、開放の電解セルでは熱の過剰を再現しなかった~\cite{nhe1998}。

マイルズらは、重水中のパラジウム陰極について、電解ガス中のヘリウム4が過剰電力とともに増したと報告した。軽水で行った対照は、ヘリウム4も過剰電力も与えなかった~\cite{miles1993}。マクキューブリらは、パラジウムの重水素吸蔵が高いときに、電力の過剰を報告した~\cite{mckubre1994}。フォカルディ、ハーベル、ピアンテッリは、ニッケルに水素ガスを吸蔵させ、熱の過剰を報告した~\cite{focardi1994}。

数キロ電子ボルトの重陽子ビームをパラジウムへ入射する入口は、別の入口である。結城、笠木らは、入射エネルギー \(2.5\) から \(10\,\mathrm{keV}\) の重陽子どうしの反応について陽子の収量を測り、より高いエネルギーからの外挿を上回る収量を報告した~\cite{yuki1998}。その運動エネルギーは、室温の運動エネルギーよりはるかに上にある。以下で計算する席は、室温の対の席である。

北村、高橋、岩村、笠木らは、その機構が2015年から2017年に行った計画で、ジルコニアに担持したパラジウムとニッケルのナノ複合体、および銅とニッケルのナノ複合体へ、重水素ガスまたは水素ガスを通した。彼らは、室温および \(200\)--\(300\,^{\circ}\mathrm{C}\) の近くで熱の過剰を報告し、分けた試料について第二の実験室が同じ過剰を報告した~\cite{kitamura2017,iwamura2017}。

戻ってくる母体は、パラジウムとニッケルである。どちらも面心立方の配列であり、八面体の隙間のまわりの最も近い正イオンは正八面体をなす。以下の節は、その正八面体にある同符号の対の、軸方向の成分を計算する。これらの報告のあいだの選択も、反応率も、熱の過剰も、導出しない。

\section{籠}
\label{sec:cage}

常温核融合は、核の同伴稿が定めた遅い荷電の入口である~\cite{dst-strong}。室温の運動エネルギーをもつ、同符号の対の入口である。ラピディティの閾値より下では、核の項はクーロン項に対抗し、閾値のすぐ下では、電磁の因子が正であるかぎり、二つの前因子のあらゆる有限な比を上回る。自由な落下は、節より前に閾値に出会う。一軸に沿う保持の相手は、電荷について既に定めた電磁チャンネルである~\cite{dst-shell}。第二の結合は導入しない。

等スケール軌跡の上で、二つの電荷の間隔を増す力は
\begin{equation}
\label{eq:force}
F(r)=\frac{k q_1 q_2}{\gamma_s(r)\,r^{2}},
\qquad
k=\frac{1}{4\pi\epsilon_0}
\end{equation}
であり、\(\gamma_s(x)=\cosh x\cos x-\sinh x\sin x\)、\(x=\ell/(2r)\) である。同符号の電荷は \(\gamma_s>0\) で斥け合い、\(\gamma_s<0\) で引き合う。第一の節 \(x_1\) は、\(\tan x=\coth x\) の \((\pi/4,\pi/2)\) における唯一の根であり、\(r_1/\ell=0.53330\) にある。外側の井戸は \(r>r_1\) であり、そこでは \(\gamma_s>0\) である。

正イオンの正八面体は、中心にさらに正のイオンを保持し、電子をその中心から押し出す。中心での力は対称性によって消え、正イオンのどの変位も戻され、電子のどの変位もさらに押し進められる。クーロンの法則のもとでは、どちらの割合も零である。その零がアーンショーの定理であり、符号つきの因子がその零を除く。外側の井戸にある金属は、保持されたイオンの配列であり、電子はそのあいだに席をもたない。遮蔽は導入しない。各イオンは点電荷のままである。

面心立方の配列の八面体の隙間のまわりで、最も近い正イオンはこの正八面体をなす。パラジウムとニッケルは、その配列のものである。この名は、隙間の幾何を記録する。第\ref{sec:reports}節の報告は、そこで核の入口が起きたかどうかを問うた。以下の席は、その配置が何を生じたかを決めない。

\section{軸方向の成分}
\label{sec:axis}

中心のイオンは電荷が一つであり、すべての方向に戻される。一軸の上の二つの陽子は対であり、ここで計算する成分は、その軸に沿う成分である。電荷 \(Ze\) のイオンを \((\pm a,0,0)\)、\((0,\pm a,0)\)、\((0,0,\pm a)\) に置き、電荷 \(+e\) の陽子を \((\pm t,0,0)\) に置く。\(0<t<a\) とする。イオンと陽子の間隔、および陽子どうしの間隔は、どの節の上にもないとする。間隔を増す成分は
\begin{equation}
\label{eq:axial}
F_{\parallel}(t)
=
-Z F(a-t)+Z F(a+t)
+4Z\frac{t}{R}F(R)
+F(2t)
\end{equation}
であり、\(R=\sqrt{a^{2}+t^{2}}\)、\(F\) は単位電荷二つに対する~\eqref{eq:force} である。第一項は、より近い軸上のイオンである。外側の井戸ではこれは斥力であり、そのイオンから陽子へ向かう向きは内を向くので、項は間隔を減らす。第二項は、より遠い軸上のイオンであり、間隔を増す。第三項は赤道上の四つのイオンである。各斥力は、外向きの軸成分 \(t/R\) をもつ。最後の項は相手の陽子である。外側の井戸では、\(F_{\parallel}\) は斥力どうしの争いである。

重陽子は同じ電荷 \(+e\) を担うので、\eqref{eq:axial} は同じ籠の中の重陽子対の軸方向の成分である。金属結合で、近傍が八つ、および十二の配位の鞍点は、減らされた電子の分量の鞍点である。それらは、この隙間の席ではない。

\section{節の外では外向き、長さの二倍では内向き}
\label{sec:seat}

籠の半径を \(a=2\ell\) に取り、半間隔を \(t=\ell\) に取る。陽子の間隔は、そのとき \(2\ell\) である。より近いイオンは \(\ell\) に立ち、より遠いイオンは \(3\ell\) に立ち、赤道上の各イオンは \(\sqrt{5}\,\ell\) に立ち、相手は \(2\ell\) に立つ。これらの距離はいずれも外側の井戸にある。イオン電荷 \(Z\ge 1\) のすべてに対して
\begin{equation}
\label{eq:inward}
F_{\parallel}(\ell)<0
\end{equation}
である。より近いイオンは、より遠いイオンと、赤道上の四つのイオンと、相手とを、あわせたものに勝つ。その位相が最大であるから、\(1/\gamma_s\) は他より強くそれを増幅する。この間隔での籠の寄与は、それ自身が内向きであり、より大きい \(Z\) はその内向きの側に加わる。相手は \(Z\) に依らない。

間隔が第一の電荷の節まで減るとき、\(2t\downarrow r_1\) であり、相手の位相は \(x_1\) まで増し、\(F(2t)\) は \(+\infty\) へ向かう。六つのイオンは節から離れたままである。最も近いイオンは \(2\ell-r_1/2\) に近づき、それはなお \(r_1\) の外側にあり、他の五つはさらに遠い。それらの寄与は有限な値へ向かうので、
\begin{equation}
\label{eq:diverge}
F_{\parallel}(t)\to +\infty
\qquad\text{（\(2t\downarrow r_1\) のとき）}
\end{equation}
である。成分は、節の右側の除かれた近傍で外向きであり、間隔 \(2\ell\) で内向きである。そのあいだの間隔の上では連続であり、それらの間隔はすべて外側の井戸に留まる。そのあいだで成分は厳密に減るので、零点は一つである。
\begin{equation}
\label{eq:balance}
r_1<s<2\ell,
\qquad
F_{\parallel}(s/2)=0.
\end{equation}
節に向かう側では成分は外向きであり、間隔 \(2\ell\) に向かう側では内向きである。どちらの側も対をその間隔へ戻す。籠は対を節の向こうへ運ばない。

接近を止める発散は、有限な半径 \(r_1\) にあり、そこでは \(\gamma_s\) が消える。クーロンの法則のもとでは、相手の項が発散するのは一致の点だけであり、一つの電荷の上の三方向の和は零であるから、籠は陽子一つを中心に保持しない。符号つきの因子は、中心に正イオン一つを保持し、対を有限な節で止める。

\section{席が開いたままにするもの}
\label{sec:open}

この席は、成分が消える唯一の間隔であり、対を \((r_1,2\ell)\) に保つ。自由な落下の相対ラピディティとは、節へ向かって落ちていく間隔のラピディティであり、席は間隔をその落下から外す。室温は、核の同伴稿の類のままである。運動エネルギーは静止エネルギーの高々 \(10^{-10}\) である。対に対する核の比較は、その入口のままである。核の節で電磁の因子は正であり、ラピディティは閾値より下である。

席は、電磁の長さにおける言明である。核の入口は、核の長さにおける言明である。籠は、その比を定めない。電荷の節が核の節の外側にあるときには、釣り合い~\eqref{eq:balance} は無限の電磁の壁の外側に立ち、軸の上の古典的な経路は核の表皮に達しない。席は、反応率も、熱の過剰も、金属の中で常温核融合が起きるかどうかを問うてきた実験のあいだの選択も、導出しない。

\begin{thebibliography}{99}

\bibitem{dst-strong}
Anonymous,
``The Strong Force as a Repulsive Layer of Proximity Gravity: Proton Chaos, Magic Numbers, Slow Neutrons, and Cold Fusion,''
companion note.

\bibitem{dst-shell}
Anonymous,
``Electron Orbits from the Signed Electrostatic Force: Covalent Bonds, Metallic Bonds, and Crystals,''
companion note.

\bibitem{fleischmann1989}
M.~Fleischmann and S.~Pons,
``Electrochemically induced nuclear fusion of deuterium,''
\textit{J.~Electroanal.~Chem.} \textbf{261}, 301 (1989).

\bibitem{jones1989}
S.~E.~Jones and others,
``Observation of cold nuclear fusion in condensed matter,''
\textit{Nature} \textbf{338}, 737 (1989).

\bibitem{yamaguchi1990}
E.~Yamaguchi and T.~Nishioka,
``Cold nuclear fusion induced by controlled out-diffusion of deuterons in palladium,''
\textit{Jpn.~J.~Appl.~Phys.} \textbf{29}, L666 (1990).

\bibitem{arata1999}
Y.~Arata and Y.-C.~Zhang,
``Observation of anomalous heat release and helium-4 production from highly deuterated palladium fine particles,''
\textit{Jpn.~J.~Appl.~Phys.} \textbf{38}, L774 (1999).

\bibitem{mizuno1996}
T.~Mizuno and others,
``Anomalous isotopic distribution of elements deposited on palladium induced by cathodic electrolysis,''
\textit{Denki Kagaku} \textbf{64}, 1160 (1996).

\bibitem{iwamura2002}
Y.~Iwamura, M.~Sakano, and T.~Itoh,
``Elemental analysis of Pd complexes: effects of D\(_{2}\) gas permeation,''
\textit{Jpn.~J.~Appl.~Phys.} \textbf{41}, 4642 (2002).

\bibitem{nhe1998}
New Energy and Industrial Technology Development Organization,
``New Hydrogen Energy Verification, Engineering and Development,''
Report NEDO-NHE-9701 (1998).

\bibitem{miles1993}
M.~H.~Miles and others,
``Correlation of excess power and helium production during D\(_{2}\)O and H\(_{2}\)O electrolysis using palladium cathodes,''
\textit{J.~Electroanal.~Chem.} \textbf{346}, 99 (1993).

\bibitem{mckubre1994}
M.~C.~H.~McKubre and others,
``Isothermal flow calorimetric investigations of the D/Pd and H/Pd systems,''
\textit{J.~Electroanal.~Chem.} \textbf{368}, 55 (1994).

\bibitem{focardi1994}
S.~Focardi, R.~Habel, and F.~Piantelli,
``Anomalous heat production in Ni-H systems,''
\textit{Nuovo Cimento A} \textbf{107}, 163 (1994).

\bibitem{yuki1998}
H.~Yuki, J.~Kasagi, and others,
``Anomalous enhancement of DD reaction in Pd and Au/Pd/PdO heterostructure targets under low-energy deuteron bombardment,''
\textit{JETP Lett.} \textbf{68}, 823 (1998).

\bibitem{kitamura2017}
A.~Kitamura and others,
``Collaborative examination on anomalous heat effect using nickel-based binary nanocomposites supported by zirconia,''
\textit{J.~Condensed Matter Nucl.~Sci.} \textbf{24}, 202 (2017).

\bibitem{iwamura2017}
Y.~Iwamura and others,
``Replication experiments at Tohoku University on anomalous heat generation using nickel-based binary nanocomposites and hydrogen isotope gas,''
\textit{J.~Condensed Matter Nucl.~Sci.} \textbf{24}, 191 (2017).

\end{thebibliography}

\end{document}
